\section*{Bab \ref{ch:b3:complex-spectra-magnetism} --- Spektrum Elektronik dan Kemagnetan Kompleks}

\begin{solution}{exo:b3:complex-spectra-magnetism:1}
D: E$_g$ + T$_{2g}$ ($2 + 3 = 5$); F: A$_{2g}$ + T$_{1g}$ + T$_{2g}$ ($1 + 3 + 3 = 7$); G: A$_{1g}$ + E$_g$ + T$_{1g}$ + T$_{2g}$
($1 + 2 + 3 + 3 = 9$).
\end{solution}

\begin{solution}{exo:b3:complex-spectra-magnetism:2}
$d^1$ $^2$D $\to$ $^2$T$_{2g}$; $d^2$ $^3$F $\to$ $^3$T$_{1g}$; $d^3$ $^4$F $\to$ $^4$A$_{2g}$; $d^4$ $^5$D $\to$ $^5$E$_g$; $d^5$ $^6$S $\to$ $^6$A$_{1g}$; $d^6$ $^5$D
$\to$ $^5$T$_{2g}$; $d^7$ $^4$F $\to$ $^4$T$_{1g}$; $d^8$ $^3$F $\to$ $^3$A$_{2g}$; $d^9$ $^2$D $\to$ $^2$E$_g$.
\end{solution}

\begin{solution}{exo:b3:complex-spectra-magnetism:3}
\ce{Mn^{2+}} spin tinggi ($n = 5$) $5.92\,\mu_B$, \ce{Fe^{2+}} (4) $4.90\,\mu_B$, \ce{Co^{2+}} (3) $3.87\,\mu_B$; \ce{Fe^{2+}} dan
\ce{Co^{3+}} spin rendah ($t_{2g}^6$): 0, diamagnetik.
\end{solution}

\begin{solution}{exo:b3:complex-spectra-magnetism:4}
$\ce{[Mn(H2O)6]^{2+}}$ (terlarang spin dan terlarang Laporte, hampir tidak berwarna) $<$ $\ce{[Co(H2O)6]^{2+}}$
(terlarang Laporte, vibronik) $<$ $\ce{[CoCl4]^{2-}}$ (tanpa pusat simetri) $<$ \ce{MnO4-} (transfer
muatan yang diizinkan).
\end{solution}

\begin{solution}{exo:b3:complex-spectra-magnetism:5}
8500: $^3$A$_{2g}\to{}^3$T$_{2g}$, jadi $\Delta_{\mathrm o} = \qty{8500}{cm^{-1}}$; 14100: $^3$T$_{1g}$(F); 24900: $^3$T$_{1g}$(P).
$15B = 39\,000 - 3 \times 8500$, $B = \qty{900}{cm^{-1}}$, $\beta = 900/1056 = 0.85$.
\end{solution}

\begin{solution}{exo:b3:complex-spectra-magnetism:6}
$\Delta_{\mathrm o} = \qty{17400}{cm^{-1}}$; dengan menyelesaikan $7.5B + 1.5\Delta_{\mathrm o} - \frac12\sqrt{225B^2 - 18B\Delta_{\mathrm o} + \Delta_{\mathrm o}^2} = 24\,600$: $B = \qty{729}{cm^{-1}}$,
$\beta = 0.79$. $\nu_3 = \qty{38500}{cm^{-1}}$ (\qty{260}{nm}), di daerah ultraviolet.
\end{solution}

\begin{solution}{exo:b3:complex-spectra-magnetism:7}
Elektron-elektron $d$ lebih terdelokalisasi ke ion-ion oksida dalam batu permata
daripada ke molekul air: ikatan Cr--O di sana lebih kovalen.
\end{solution}

\begin{solution}{exo:b3:complex-spectra-magnetism:8}
Kuat ($e_g$ terisi tidak merata): $d^9$, $d^4$ spin tinggi, $d^7$ spin rendah. Tidak ada: $d^3$ ($t_{2g}^3$, suku A) dan
$d^6$ spin rendah ($t_{2g}^6$). Lemah: $d^6$ spin tinggi ($t_{2g}^4e_g^2$, suku T).
\end{solution}

\begin{solution}{exo:b3:complex-spectra-magnetism:9}
$d^7$ oktahedral spin tinggi: suku dasar $^4$T$_{1g}$, suku T dengan \omterm{def:b3:complex-spectra-magnetism:moment}{kontribusi orbital},
momennya jauh di atas nilai spin saja $3.87\,\mu_B$. $d^7$ tetrahedral ($\equiv$ $d^3$ oktahedral): suku
dasar $^4$A$_2$, momen orbitalnya terpadamkan, dekat dengan nilai spin saja (sedikit di atasnya,
karena tercampur dengan suku-suku T tereksitasi).
\end{solution}

\begin{solution}{exo:b3:complex-spectra-magnetism:10}
$\mu_{\mathrm{eff}} = 2.828\sqrt{1.87} = 3.87\,\mu_B$: tiga elektron yang tidak berpasangan ($S = \frac32$).
\end{solution}

\begin{solution}{exo:b3:complex-spectra-magnetism:11}
$\Delta\chi_v = 3 \times 25.0/\num{400e6} = \num{1.88e-7}$; $\chi_{\mathrm m} = \num{1.88e-7}/\qty{10.0}{mol.m^{-3}} = \qty{1.88e-8}{m^3.mol^{-1}}$, yaitu
\qty{1.49e-3}{cm^3.mol^{-1}}; $\chi_{\mathrm m}T = 0.445$, $\mu_{\mathrm{eff}} = 1.89\,\mu_B$: satu elektron yang tidak berpasangan.
\end{solution}

\begin{solution}{exo:b3:complex-spectra-magnetism:12}
$T_{1/2} = 15\,000/75 = \qty{200}{K}$. Pada \qty{250}{K}: $\Delta G = 15\,000 - 250 \times 75 = \qty{-3750}{J/mol}$, $K = \eu^{1.80} = 6.07$,
$x_{\mathrm{HS}} = 0.86$; $\chi_{\mathrm m}T = 0.86 \times 3.00 = \qty{2.6}{cm^3.K.mol^{-1}}$.
\end{solution}

\begin{solution}{pb:b3:complex-spectra-magnetism:1}
\textbf{1.} [Ar]$3d^3$.
\textbf{2.} $^4$F.
\textbf{3.} $^4$A$_{2g}$ + $^4$T$_{2g}$ + $^4$T$_{1g}$.
\textbf{4.} $^4$A$_{2g}$, suku $t_{2g}^3$, dengan ketiga elektron di orbital-orbital yang lebih rendah.
\textbf{5.} $^4$T$_{1g}$(P).
\textbf{6.} $^4$A$_{2g}\to{}^4$T$_{2g}$, $\to{}^4$T$_{1g}$(F), $\to{}^4$T$_{1g}$(P).
\textbf{7.} 561 dan \qty{414}{nm}.
\textbf{8.} 597 dan \qty{434}{nm}.
\textbf{9.} Rubi menyerap kuning-hijau dan ungu serta meneruskan merah (dengan
sedikit biru); pita-pita zamrud, yang bergeser ke arah merah, juga menyerap
jingga-merah, dan hijaulah yang diteruskan.
\textbf{10.} Rubi \qty{17820}{cm^{-1}}, zamrud \qty{16740}{cm^{-1}}.
\textbf{11.} $^4$T$_{2g}$ ($t_{2g}^2e_g$) dan $^4$A$_{2g}$ ($t_{2g}^3$) memiliki energi tolakan elektron yang sama; selisih
keduanya adalah energi promosi satu-elektron $\Delta_{\mathrm o}$.
\textbf{12.} $B = (14\,305 - 547)/15 = \qty{917}{cm^{-1}}$.
\textbf{13.} \qty{614}{cm^{-1}}.
\textbf{14.} \qty{618}{cm^{-1}}.
\textbf{15.} 0.67 untuk keduanya.
\textbf{16.} 29.0 dan 27.1.
\textbf{17.} $\nu_3 = \qty{38500}{cm^{-1}}$ (\qty{260}{nm}) untuk rubi, di daerah ultraviolet, tempat serapan
transfer muatan menyembunyikannya.
\textbf{18.} Situs oktahedral beril lebih besar: ikatan Cr--O yang lebih panjang memberikan
pembelahan yang lebih kecil, yang berubah tajam bersama jarak.
\textbf{19.} $10^7/695 = \qty{14390}{cm^{-1}}$, $23.4B$.
\textbf{20.} $^2$E$_g$ terutama bergantung pada $C$: tingkat terukur yang lebih tinggi berarti $C/B$ lebih besar,
sekitar 5.3 menurut diagonalisasi yang sama.
\textbf{21.} $^2$E$_g$ dan $^4$A$_{2g}$ sama-sama berasal dari $t_{2g}^3$ dan hanya berbeda oleh pembalikan sebuah spin: ikatannya,
jadi juga geometrinya, sama pada keduanya, dan transisinya tidak memiliki
progresi vibrasi (sebuah garis 0--0).
\textbf{22.} Emisi itu terlarang spin.
\textbf{23.} Energinya bergantung pada $B$ dan $C$, hampir tidak pada $\Delta_{\mathrm o}$ (garis datar pada diagram),
dan $B$ hampir sama pada kedua permata.
\textbf{24.} Ion-ion yang dipompa ke dalam pita-pita lebar berelaksasi dan menumpuk di tingkat
$^2$E$_g$ yang berumur panjang, sehingga jumlah ion di tingkat itu dapat melebihi jumlah ion di tingkat dasar:
sebuah inversi \omterm{def:b3:partition-functions:boltzmann}{populasi}.
\textbf{25.} \textbf{$\Delta_{\mathrm o}(\text{rubi}) - \Delta_{\mathrm o}(\text{zamrud}) = \qty{1080}{cm^{-1}}$}: seluruh perbedaan antara merah dan
hijau.
\end{solution}
